\section*{Введение}
\addcontentsline{toc}{section}{Введение}

\textbf{Актуальность темы.} На современном этапе глобального исторического развития, характеризующемся переходом к шестому технологическому укладу и тотальной цифровизацией социально-экономических подсистем, научно-техническая интеллигенция (НТИ) становится ключевым актором макрополитических и цивилизационных процессов. В условиях беспрецедентного геополитического противостояния, жесткого санкционного давления и необходимости форсированного достижения технологического суверенитета Российской Федерации, интеллектуальный потенциал инженеров, ученых, конструкторов и IT-специалистов превращается в главный стратегический ресурс выживания и развития государства. 

Историческая ретроспектива показывает, что взаимоотношения между властью и носителями научного знания в России всегда носили амбивалентный и глубоко драматичный характер. Начиная с петровских модернизационных сдвигов, через советский феномен «сверхиндустриализации» и вплоть до современных реалий цифровой экономики, государство выступало главным заказчиком, патроном и одновременно контролером научно-технического сословия. В то же время сама интеллигенция перманентно балансировала между статусом элитарной «техноструктуры», непосредственно влияющей на государственные решения, и положением интеллектуального пролетариата, отчужденного от реальных рычагов управления. Исследование социологического портрета, ценностных ориентаций, механизмов лоббирования и протестного потенциала НТИ через призму теории политических элит и групп интересов позволяет эксплицировать латентные пружины модернизации и спрогнозировать устойчивость государственного каркаса, что определяет высокую научно-теоретическую и прикладную значимость данной работы.

\textbf{Объект исследования} --- научно-техническая интеллигенция как специфическая социально-профессиональная группа и субъект модернизационных процессов в исторической динамике российского общества.

\textbf{Предмет исследования} --- структурно-функциональные особенности, политико-правовой статус, социологические характеристики и механизмы взаимодействия научно-технической интеллигенции с институтами государственной власти в современной России.

\textbf{Цель работы} --- осуществить комплексный историко-политологический и социологический анализ эволюции, современного состояния, кризисных факторов и векторов трансформации научно-технической интеллигенции в системе властных отношений современного российского государства.

Для достижения поставленной цели необходимо решить следующие \textbf{задачи}:
\begin{enumerate}
    \item Исследовать теоретико-методологические подходы к анализу интеллигенции, концепцию «техноструктуры» Дж. Гэлбрейта и механизмы артикуляции интересов в неокорпоративистских моделях.
    \item Проследить историческую эволюцию статуса научно-технических кадров в России от истоков государственности до распада советской системы.
    \item Верифицировать социологический портрет, ценностные ориентации и стратификационные сдвиги внутри современной российской НТИ.
    \item Проанализировать формы политического участия НТИ, феномен «утечки мозгов», цифровую ресоциализацию и выработать рекомендации по оптимизации государственной научно-технической политики.
\end{enumerate}